\clearpage
\begin{summarybox}

\textbf{\textcolor{CSummary}{第一部分：$0/0$ 型——它解决什么问题？}}

当分子、分母都越来越接近 $0$ 时，仅凭代入无法判断比值；
这时可以比较它们的\textbf{变化快慢}，在满足条件时用分别求导来简化计算。

\medskip
\textbf{\textcolor{CSummary}{一句话记住}}

\[
\boxed{\displaystyle
\lim_{x\to a}\frac{f(x)}{g(x)}=\lim_{x\to a}\frac{f'(x)}{g'(x)}
\quad\text{（须先检查使用条件）}.}
\]
注意：求的是\textbf{分子导数比上分母导数}，不是整个分式的导数。

\medskip
\textbf{\textcolor{CSummary}{三步判断}}
\begin{enumerate}[leftmargin=*, itemsep=0.3em, topsep=0.3em]
  \item \textbf{先代入看情况：}是不是分子、分母都趋于 $0$？
  \item \textbf{再检查能否求导：}附近能求导、分母导数不为 $0$，并且求导后的比值有确定的极限。
  \item \textbf{最后求值：}约分或直接代入；若还满足同样的情况，才考虑再求一次导。
\end{enumerate}

\medskip
\textbf{\textcolor{CSummary}{优先顺序}}

\textbf{直接代入 $\rightarrow$ 因式分解或有理化 $\rightarrow$ 已知基本极限 $\rightarrow$ 条件满足时考虑分别求导。}

\medskip
\textbf{\textcolor{CSummary}{三个值得记住的结果}}
\[
\lim_{x\to a}\frac{x^2-a^2}{x-a}=2a,
\qquad
\lim_{x\to0}\frac{\sin x}{x}=1,
\qquad
\lim_{x\to0}\frac{x-\sin x}{x^3}=\frac16.
\]

\textbf{第一部分证明范围：}导数定义、共同零因子的多项式、单位圆比较，严格覆盖若干特定情况；不是一般法则的全部证明。
求导后的比值没有极限，也\textbf{不能}据此断言原式没有极限。

\end{summarybox}

\clearpage
\begin{summarybox}
\textbf{\textcolor{CSummary}{第二部分：$\infty/\infty$ 型——它解决什么问题？}}

当分子、分母的\textbf{绝对值}都无限增大时，只知道“两边都大”仍无法计算比值；答案由\textbf{谁增长得更快}决定。
\[
\boxed{\displaystyle
\lim_{x\to b}\frac{f(x)}{g(x)}=\lim_{x\to b}\frac{f'(x)}{g'(x)}
\quad\text{（同一方向、使用条件满足时）}.}
\]
这里 $b$ 可以表示有限位置（可取单侧）或 $+\infty$、$-\infty$；两个极限必须沿\textbf{题目指定的同一方向}计算。

\medskip
\textbf{\textcolor{CSummary}{三步判断}}
\begin{enumerate}[leftmargin=*, itemsep=0.3em, topsep=0.3em]
\item \textbf{先看类型：} $|f(x)|$、$|g(x)|$ 是否都无限增大？（可能是趋向正无穷，也可能是负无穷。）
\item \textbf{再验条件：} 附近都能求导、$g'(x)\ne0$，而且导数之比有确定的极限（含确定的正、负无穷）。
\item \textbf{最后求值：} 分别求导后能算就结束；若想继续用，需要重新检查新式子。
\end{enumerate}

\medskip
\textbf{\textcolor{CSummary}{三个典型答案}}
\[
\lim_{x\to+\infty}\frac{x}{x^2}=0,\qquad
\lim_{x\to+\infty}\frac{3x^2}{2x^2}=\frac32,\qquad
\lim_{x\to+\infty}\frac{x^3}{x^2}=+\infty.
\]

\medskip
\textbf{\textcolor{CSummary}{眼前一亮：两个小提醒}}

\textbf{一是：} $\infty/\infty$ 型不要求 $x$ 一定趋向无穷大；$x$ 接近有限位置时，函数值也可能都趋向无穷大。

\textbf{二是：} $x\to-\infty$ 时，$\sqrt{x^2}=|x|=-x$，不能直接写成 $x$。

\medskip
\textbf{\textcolor{CSummary}{高中最省时的方法与证明范围}}

\textbf{先比较最高次项 $\rightarrow$ 再约分、化根号 $\rightarrow$ 不好直接化简且满足条件时考虑分别求导。}

本部分\textbf{严格验证}了多项式分式，以及具体的根式、对数式；不宣称它们就是一般法则的完整证明。即便导数比值不存在，原式的极限也可能存在。
\end{summarybox}
